\documentclass[11pt,a4paper]{article}
\usepackage[T1]{fontenc}
\usepackage{lmodern,amsmath,amssymb,booktabs,tabularx,microtype}
\usepackage[margin=25mm]{geometry}
\usepackage{hyperref}
\hypersetup{colorlinks=true,linkcolor=blue,urlcolor=blue,
pdftitle={Retrieval, Association, and Shrinkage in BundleFuse}}
\setlength{\parindent}{0pt}
\setlength{\parskip}{5pt}
\setlength{\emergencystretch}{3em}
\newcommand{\pos}[1]{\left[#1\right]_+}
\newcommand{\diag}{\operatorname{diag}}
\newcommand{\off}{\operatorname{off}}
\title{Retrieval, Association, and Shrinkage in BundleFuse\\
\large What the connections to established methods actually explain}
\author{}\date{}
\begin{document}
\maketitle

The connection to text retrieval begins with a simple observation: a historical
bundle is a collection of items, just as a document is a collection of terms.
An observed partial bundle can therefore be matched against historical bundles
using occurrence, frequency and size information. Completion then requires an
additional step: the matched bundles transfer their support to candidate items.
The equations below make this interpretation exact for the binary, unit-weight
training bundles used in the full-ranking POG formulation.

\section{Define the observations and the retrieval analogy}
Let $\mathcal I=\{1,\ldots,n\}$ be the item catalogue, and let $m$ training
bundles be indexed by $e\in\{1,\ldots,m\}$. Define the binary incidence matrix
$B\in\{0,1\}^{m\times n}$ by
\[
 B_{ei}=\begin{cases}1,&\text{item $i$ occurs in training bundle $e$},\\
 0,&\text{otherwise}.
 \end{cases}
\]
An item contributes at most once to a training row. Define its training frequency
$f_i=\sum_{e=1}^mB_{ei}$ and bundle size $d_e=\sum_{i=1}^nB_{ei}$.
Let $S\subseteq\mathcal I$ be a nonempty observed set, let $c\in\mathcal I\setminus S$
be a candidate, and let $x_S\in\{0,1\}^n$ be the indicator of $S$.

\begin{center}
\begin{tabularx}{\linewidth}{@{}XX@{}}\toprule
Text retrieval & Bundle completion \\\midrule
Document & Historical training bundle $e$ \\
Term & Catalogue item $i$ \\
Occurrence of a term in a document & Binary entry $B_{ei}$ \\
Number of documents containing a term & Number of bundles containing an item, $f_i$ \\
Document length & Bundle size $d_e$ \\
Query terms & Observed items $S$ \\\bottomrule
\end{tabularx}
\end{center}

In a basic TF--IDF formulation, term $i$ in document $e$ receives weight
\begin{equation}
 \operatorname{tfidf}(i,e)=\operatorname{tf}(i,e)
                         \operatorname{idf}(i),\qquad
 \operatorname{idf}(i)=\log\frac{m}{f_i}\quad(f_i>0).
\end{equation}
Here $\operatorname{tf}$ counts occurrences within the document, while
$\operatorname{idf}$ discounts terms appearing in many documents. Common variants
smooth the logarithm and normalize document vectors by their Euclidean length
\cite{tfidf}. With binary bundle rows, the corresponding occurrence quantity is
simply $B_{ei}$; repeated occurrences within a row are not modeled.

This analogy motivates the roles of the weights. It does not determine their
exact functional form. BundleFuse uses inverse square roots and explicit offsets,
as defined next, rather than the logarithmic IDF above.

\section{Match the observed items to historical bundles}
Define an observed-item weight $a_i$ and a historical-bundle weight $v_e$ by
\begin{equation}
 a_i=\frac1{\sqrt{f_i+5}},\qquad
 v_e=\frac1{\sqrt{d_e+1}}.
\end{equation}
The historical bundle's match score is
\begin{equation}\label{overlap}
 \boxed{o_e(S)=v_e\sum_{s\in S}B_{es}a_s
 =\frac1{\sqrt{d_e+1}}
   \sum_{s\in S}\frac{B_{es}}{\sqrt{f_s+5}}.}
\end{equation}
Each summand corresponds to one observed item actually present in the historical
bundle. Its frequency discount is larger for a rarer observed item. The outer
factor reduces the contribution of the same match when it occurs inside a larger
bundle. Multiple matches add together.

For example, an observed item appearing in four training bundles has weight
$1/\sqrt9=1/3$, whereas one appearing in 95 bundles has weight
$1/\sqrt{100}=1/10$. Within the same historical bundle, the former match
contributes $10/3$ times as much. Likewise, the bundle-size factors for sizes
three and fifteen are $1/2$ and $1/4$, respectively: the same weighted overlap
receives twice as much weight in the smaller bundle. These are illustrative
calculations, not observations about a particular evaluation query.

\paragraph{What is shared with TF--IDF?}
Both constructions discount a clue according to how widely it occurs in the
training collection. Document-vector normalization can additionally reduce length effects;
IDF alone does not normalize document length. This is a parallel in weighting principles, with
explicit differences in the formulas:
\begin{itemize}
\item The item-frequency factor $(f_i+5)^{-1/2}$ is a smoothed decreasing
frequency weight, rather than logarithmic IDF.
\item The factor $(d_e+1)^{-1/2}$ depends on binary bundle size. It is not the
Euclidean normalization of a TF--IDF-weighted document vector, whose norm would
also depend on its individual term weights.
\item Frequency weights act on the observed query items in~\eqref{overlap}.
Candidate items do not receive that same divisor when votes are transferred.
\end{itemize}
The offsets $5$ and $1$ and the square-root exponents are fixed modeling
choices. The analogy explains their roles, not the optimality of those numbers.
There is a further distinction: if $f_i=m$, basic unsmoothed IDF is zero,
whereas $a_i=(m+5)^{-1/2}$ remains positive. Nor is $a_i$ a logarithmic
information content. Thus the direct score is best described as
\emph{frequency-discounted historical retrieval}, not as TF--IDF itself.
These are discounts per recorded match; repeated matches can still accumulate
large total support. They do not establish complete removal of popularity effects.

\section{Transfer historical-bundle matches to completion candidates}
A historical bundle proposes its member items as possible completions.
Accordingly, sum its match score into candidate $c$ only if it contains $c$:
\begin{equation}\label{pair}
 \boxed{p_c(S)=\sum_{e=1}^mB_{ec}o_e(S)
 =\sum_{s\in S}a_s\sum_{e=1}^m v_eB_{ec}B_{es}.}
\end{equation}
The inner sum is the size-discounted number of bundles containing both the
candidate and observed item $s$. The outer sum combines this support across the
observed items, with their individual frequency discounts.

Set $D=\diag(a_i)\in\mathbb R^{n\times n}$ and
$V=\diag(v_e)\in\mathbb R^{m\times m}$. The complete calculation is
\begin{equation}
 o(S)=VBDx_S,\qquad \boxed{p(S)=B^\top VBDx_S.}
\end{equation}
Reading from right to left: weight the observed items, match them to historical
bundles, discount bundle size, and transfer the bundle scores to their members.
This provides a retrieval interpretation of the two incidence multiplications.
The first multiplication produces historical-bundle evidence; the second
converts that evidence into a completion ranking.

The formula retains self-contributions at observed-item entries of $p$.
Candidates are restricted to $c\notin S$ when recommendations are returned.
The displayed equation is the raw direct score, before catalogue standardization
or combination with other sources.

\section{Distinguish simultaneous matches from separate matches}
Additive overlap gives one contribution for each observed item. To reward several
observed items appearing in the same historical bundle, define
\[
 k_e(S)=\sum_{s\in S}B_{es}=|e\cap S|,\qquad [z]_+=\max(z,0),
 \qquad\beta=1.5.
\]
The incidence-enhanced score is
\begin{equation}\label{group}
 \boxed{g_c(S)=\sum_eB_{ec}o_e(S)
                   \bigl[1+\beta[k_e(S)-1]_+\bigr].}
\end{equation}
For one matching observed item the multiplier is one. For two it is $2.5$;
for three it is four. If there are no matches, $o_e(S)=0$.
This amplification applies separately to each historical bundle before its
vote is transferred. The setting $\beta=1.5$ specifies the chosen strength.

Why retain this operation? A candidate may occur with $s$ in one historical
bundle and with $t$ in another, without ever occurring with both together.
Equation~\eqref{group} distinguishes that situation from a historical bundle
containing $s,t,c$. To see the distinction algebraically, define
\[
 T_{cst}=\sum_e v_eB_{ec}B_{es}B_{et}.
\]
This is the weighted count of bundles containing all three distinct items.
For $c\notin S$, the binary-query expansion of~\eqref{group} is
\begin{equation}
 g_c(S)=p_c(S)+\beta\sum_{\{s,t\}\subseteq S}(a_s+a_t)T_{cst},
\end{equation}
where each unordered pair of distinct observed items appears once.
Thus direct matching uses candidate--seed pairs, while the correction uses
candidate--seed--seed triples. It rewards supported pairs within $S$; it does
not require the entire observed set to have appeared together.

The TF--IDF analogy explains the frequency and size discounts. The simultaneous
term extends the scoring rule beyond additive item matching. Its triple counts
cannot in general be reconstructed from the item-pair count matrix alone.
This is a statement about information retained by these particular operators,
not a claim that multi-item retrieval or higher-order recommendation is new.

\section{Interpret PPMI as conditional occurrence plus inverse-frequency information}
The PPMI source uses another normalization of the same training observations.
Define
\begin{equation}
 C=\off(B^\top B),\qquad
 q_i=\sum_jC_{ij},\qquad T=\sum_{i,j}C_{ij},
\end{equation}
where $\off$ replaces a square matrix's diagonal by zero. $C_{ij}$ counts
training bundles containing both distinct items $i,j$. Assume $T>0$.
Sampling an ordered within-bundle pair from all recorded pair occurrences gives
\[
 P(i,j)=\frac{C_{ij}}{T},\qquad
 P(i)=\frac{q_i}{T}.
\]
The endpoint marginal $q_i$ is not the bundle frequency $f_i$:
\begin{equation}
 q_i=\sum_eB_{ei}(d_e-1).
\end{equation}
An item in a larger bundle participates in more pairs. This distinction must be
retained when connecting the log-ratio to inverse frequency.

For a recorded pair $(c,s)$, its pointwise mutual information is
\begin{align}
 \operatorname{PMI}(c,s)
 &=\log\frac{P(c,s)}{P(c)P(s)}
   =\log\frac{C_{cs}T}{q_cq_s}\\
 &=\underbrace{\log\frac{C_{cs}}{q_c}}_{\log P(s\mid c)}
  +\underbrace{\log\frac{T}{q_s}}_{-\log P(s)}.\label{pmi}
\end{align}
The first term measures how often $s$ appears as a partner of $c$; the second
is the inverse-frequency information of $s$ in the ordered-pair distribution.
The sum is positive exactly when $s$ is more likely beside $c$ than under its
overall endpoint frequency. This is the precise information-theoretic parallel
to discounting common terms.

Define positive PMI by $M_{cs}=[\operatorname{PMI}(c,s)]_+$ for recorded pairs,
and zero for absent pairs. The raw source score is
\begin{equation}
 r_{\rm pmi}(c\mid S)=\sum_{s\in S}M_{cs}.
\end{equation}
The implementation adds small numerical stabilizers to the denominators and
logarithm. Equation~\eqref{pmi} describes the underlying ratio before those
stabilizers. Clipping removes negative associations from this source; it does
not convert them into penalties. Unlike TF--IDF, PPMI includes the conditional
association of two items, rather than only a term's frequency weight.


\section{From association to similarity of association profiles}
The direct PPMI entry $M_{cs}$ asks whether $c$ and $s$ occur together above
the independent-endpoint baseline. A different question is whether they have
similar partners. This distinction explains the SVD source.

Write $M=U\Sigma V^\top$ for a singular value decomposition, with singular
values in decreasing order. Let $r$ be the retained dimension, and $V_r$ the
first $r$ right singular vectors. Define
\[
 Z=MV_r=U_r\Sigma_r,\qquad
 u_i=\frac{Z_i}{\|Z_i\|_2+10^{-8}},\qquad
 r_{\rm svd}(c\mid S)=u_c^\top\frac1{|S|}\sum_{s\in S}u_s.
\]
Here $Z_i$ is the row for item $i$, not a singular vector. Orthogonality of the
columns of $V_r$ does not make these rows unit length.
Before row normalization,
\begin{equation}\label{profile}
 ZZ^\top=MV_rV_r^\top M^\top=U_r\Sigma_r^2U_r^\top.
\end{equation}
If all singular directions are retained, its $(c,s)$ entry is
$\sum_j M_{cj}M_{sj}$. Thus it rewards sharing positively associated partners,
even when $M_{cs}=0$. Truncation retains the dominant directions in these
association profiles. It is not the rank-$r$ reconstruction
$U_r\Sigma_rV_r^\top$ used directly as an item--item score.

The connection is to the use of SVD in latent semantic indexing
\cite{lsi}: compare objects through compressed occurrence profiles.
Here the input is item--item PPMI, not a term--document table.
The exact rank-$r$ SVD minimizes $\|M-A\|_F^2$ over matrices $A$ of rank at
most $r$, where $\|\cdot\|_F^2$ sums squared entries. The numerical procedure
approximates that decomposition. Neither this approximation criterion nor
Equation~\eqref{profile} guarantees better recommendation rankings.

This also limits the interpretation: similar partners may indicate that two
items are substitutes. Their similarity is not proof that they belong together.
The source supplements observed co-membership rather than replacing its meaning.

\section{From descriptions to deviations from a category mean}
Association profiles require interaction observations. Supplied item descriptions
provide another way to compare items. For one description modality, let
$X\in\mathbb R^{n\times p}$ contain its row-normalized supplied vectors,
with $p$ feature coordinates. Let $\ell(i)$ be the supplied category label,
and let $H\in\{0,1\}^{n\times k}$ contain one indicator column for each of the
$k$ nonempty categories: $H_{ia}=\mathbf1[\ell(i)=a]$.

Define
\begin{equation}
 \Pi=H(H^\top H)^{-1}H^\top,\qquad R=(I_n-\Pi)X.
\end{equation}
$I_n$ is the $n\times n$ identity matrix. Since $H^\top H$ contains the
category sizes on its diagonal, $(\Pi X)_i$ is precisely the mean vector
of all catalogue items in category $\ell(i)$.
Moreover,
\[
 \widehat A=\arg\min_{A\in\mathbb R^{k\times p}}\|X-HA\|_F^2
           =(H^\top H)^{-1}H^\top X,\qquad H^\top R=0.
\]
Thus category subtraction is exactly the residual of a least-squares fit using
category indicators. It removes the part represented by a category-specific
constant vector. This is a more precise connection than calling it
``disentanglement.''

Normalize each residual row, $E_i=R_i/(\|R_i\|_2+10^{-8})$. The descriptor
score is
\[
 r_E(c\mid S)=E_c^\top\bar E_S,\qquad
 \bar E_S=\frac1{|S|}\sum_{s\in S}E_s.
\]
It compares the candidate to the mean description of the observed items.
This is the same centroid-matching operation that underlies the positive
feedback part of Rocchio retrieval \cite{rocchio}; it is not the full Rocchio
algorithm with an original query and negative feedback.

Row normalization generally destroys $H^\top E=0$, even though $H^\top R=0$.
Consequently, only the unnormalized residual has the exact orthogonality
property. Removing a category mean does not remove all category information
or establish that remaining coordinates represent compatibility.
For several modalities the operation is performed separately, followed by
concatenation and normalization where required.

\section{From descriptions of an item to descriptions of its partners}
Two items can have related roles even if their own descriptions differ.
The bridge incorporates the descriptions of historically connected items.
Let $E$ now denote the normalized concatenation of the available residual
modalities. With $D$ already defined as the seed-frequency diagonal, set
\[
 L=\off\!\left(D B^\top
       \diag\bigl((d_e+1)^{-1}\bigr) B D\right),\qquad
 N=\mathcal N(LE),\qquad F=\mathcal N(E+\eta N),\quad\eta=0.5,
\]
where $\mathcal N$ divides each row by its Euclidean norm plus $10^{-8}$.
In particular, $(LE)_i=\sum_jL_{ij}E_j$ aggregates the residual descriptions
of items sharing training bundles with $i$. $L$ is symmetric but is not a
row-stochastic transition matrix: its rows need not sum to one.

The bridge score is $r_F(c\mid S)=F_c^\top |S|^{-1}\sum_{s\in S}F_s$.
Let $A$ be the diagonal of the final row-normalization factors, so that
$F=A(E+\eta N)$. Its similarity matrix expands exactly as
\begin{equation}
 FF^\top=A\bigl(EE^\top+\eta EN^\top+\eta NE^\top+
                   \eta^2 NN^\top\bigr)A.
\end{equation}
The four terms compare own descriptions, own descriptions with partner
profiles in either direction, and partner profiles with one another.
This is a one-step graph aggregation followed by similarity matching.
It provides a concrete route from descriptive resemblance to historical
context; it is neither an iterated random walk nor a fitted graph neural network.
The cross terms can be negative because category residual coordinates are signed.

\section{From exact identities to pooled category or cluster relations}
Rare items may have too few observations to estimate useful item-specific
associations. Grouping items shares observations across their labels or
feature memberships. This is pooling of counts, followed by lifting the
resulting relation back to individual items.

For the category indicator $H$ defined above, $BH$ counts category occurrences
per bundle. Write
\[
 A^{\rm cat}=BH,\quad C^{\rm cat}=\off((A^{\rm cat})^\top A^{\rm cat}),\quad
 h_a=\sum_e A^{\rm cat}_{ea},\quad h_+=\sum_a h_a.
\]
For positive recorded counts define
$Q^{\rm cat}_{ab}=[\log(C^{\rm cat}_{ab}h_+/(h_a h_b))]_+$, and use zero
otherwise. Small numerical stabilizers are omitted here.
The category association score is $(H Q^{\rm cat} H^\top x_S)_c$.
The matrix products explicitly mean: collect seed labels, score relations
between labels, then give each candidate the score of its label.
These $h_a$ are category-occurrence totals, not the endpoint totals of
$C^{\rm cat}$. This is a count-normalized log score; the exact independence
interpretation established for item PPMI must not be transferred to it.

The POG category score also discourages categories that rarely repeat in one
bundle. Let $n_a^{(1)}$ and $n_a^{(2)}$ count training bundles containing at
least one and at least two items of category $a$, respectively. Set
\[
 \rho_a=\frac{n_a^{(2)}+0.5}{n_a^{(1)}+1}.
\]
Subtract $0.5(1-\rho_a)$ for each seed of the candidate's category if
$\rho_a<0.15$. Under a Bernoulli model of repetition conditional on category
presence, $\rho_a$ is exactly the posterior mean with a
$\operatorname{Beta}(1/2,1/2)$ prior. The threshold and penalty magnitude
are additional design choices, not consequences of that posterior mean.

Soft groups replace $H$ by a matrix $J\in\mathbb R^{n\times k}$ of
nonnegative feature-cluster memberships. Clusters are fitted by minibatch
$k$-means to normalized concatenated supplied descriptors; membership is
proportional to the exponential of similarity to a normalized cluster centre,
with temperature $0.05$. Each row sums approximately to one (the normalization
includes a small stabilizer). For binary training bundles,
\[
 C^{\rm soft}=\off(J^\top C J),\qquad b_a=\sum_iJ_{ia},\quad b_+=\sum_a b_a,
\]
\[
 Q^{\rm soft}_{ab}=
 \left[\log\frac{C^{\rm soft}_{ab}b_+}{b_a b_b}\right]_+,
 \qquad r_{\rm soft}=JQ^{\rm soft}J^\top x_S,
\]
using zero for absent counts and the same numerical qualifications.
The identity $J^\top C J=(BJ)^\top BJ-J^\top\diag(f)J$ explains the
self-item subtraction: the same item must not supply both ends of a pair.
Here $b_a$ is catalogue membership mass, not a pair marginal.

This is relation sharing through groups. It is not a dot-product kernel in
general: $Q^{\rm soft}$ need not be positive semidefinite. Given $J$, the
soft count matrix contains no information beyond the item-pair matrix $C$.
This contrasts with the incidence correction, which uses candidate--seed--seed
counts from the original bundles. The selected full POG\_DENSE recipe does
not use the soft source; its role is to explain an available component, not
an observed cause of the dense result.

\section{From overlapping clues to diversity of support}
The optional agreement score asks whether a candidate is supported by different
observed items rather than nearly duplicate clues. For POG, define
\[
 W=\off(B^\top VBD),\qquad
 q_{cs}=W_{cs}+0.15[E_c^t\!\cdot E_s^t]_+,\qquad A_c=\sum_{s\in S}q_{cs},
\]
where $E_i^t$ is the normalized category-residual title vector defined above.
For distinct $s,t\in S$, set
\[
 K_{st}=\left[1-0.5\mathbf1[\ell(s)=\ell(t)]
                  -0.5[E_s^t\!\cdot E_t^t]_+\right]_+,\qquad K_{ss}=0.
\]
For $A_c>0$, let $\pi_{cs}=q_{cs}/A_c$. These are nonnegative weights summing
to one over seeds. Then
\begin{equation}
 r_R(c\mid S)=\frac{\sum_{s,t}q_{cs}K_{st}q_{ct}}{A_c+0.05}
 =\frac{A_c^2}{A_c+0.05}\sum_{s,t}\pi_{cs}\pi_{ct}K_{st}.
\end{equation}
When $A_c=0$ the score is zero. The last sum is the expected dissimilarity
of two independent seed draws from the candidate's support distribution.
It has the quadratic-diversity form introduced by Rao \cite{rao}.
It is multiplied here by a support-strength factor. A candidate supported by
one seed has zero diversity; strong support spread across dissimilar seeds
can receive more credit. ``Independent draws'' describes this mathematical
sampling operation, not statistical independence of the original observations.
No triangle inequality or entropy-concavity property of $K$ is asserted.

This is not the incidence correction: it combines candidate--seed affinities
and seed dissimilarities, without requiring a historical bundle containing
all three items. It is also not a penalty for diversity of the recommended
list. The selected full POG\_DENSE recipe gives this source zero weight.

\section{From heterogeneous scores to a regularized final estimate}
The sources answer different questions and have different scales.
For each raw source $r_h(i\mid S)$, define, over all $n$ catalogue items,
\[
 \mu_h=\frac1n\sum_i r_h(i\mid S),\qquad
 \sigma_h^2=\frac1{n-1}\sum_i(r_h(i\mid S)-\mu_h)^2,\qquad
 z_h(i\mid S)=\frac{r_h(i\mid S)-\mu_h}{\sigma_h+10^{-9}}.
\]
This is the full-ranking POG standardization, applied separately per query.
It expresses scores relative to their catalogue mean and spread; it is not
probability calibration. No held-out completion label enters these statistics.
The full catalogue, including seeds before the final exclusion mask, is used.

Nonnegative scalar weights combine these standardized sources. They are chosen
on validation data and need not sum to one. Therefore source fusion is not
itself necessarily a convex average. For the selected full POG\_DENSE recipe,
write
\begin{align}
 Y_c&=0.5z_{\rm cat}+0.25z_{\rm bridge}+0.5z_g
          +0.25z_{\rm pop}+0.25z_{\rm svd},\\
 P_c&=0.5z_{\rm cat}+0.25z_{\rm bridge}+0.25z_{\rm pop},
 \qquad r_{\rm pop}(c)=\log(1+f_c),
\end{align}
where every $z$ on the right is evaluated at $(c\mid S)$.
The popularity source provides a query-independent frequency baseline before
fusion, with logarithmic compression of large counts. Its increasing frequency
weight has a different role from the decreasing weight on observed seed clues.
$P_c$ is a restricted fallback score, not an external probability distribution.

Let $s_0\geq0$ be its validation-selected scale and $\kappa>0$ the frequency
threshold parameter. The final score is
\begin{equation}\label{shrink}
 \boxed{\theta(c\mid S)=\frac{f_cY_c+\kappa s_0P_c}{f_c+\kappa}.}
\end{equation}
This is exactly the unique solution of
\begin{equation}
 \min_{t\in\mathbb R}\quad
 \frac{f_c}{2}(t-Y_c)^2+\frac{\kappa}{2}(t-s_0P_c)^2.
\end{equation}
Differentiating gives $(f_c+\kappa)t=f_cY_c+\kappa s_0P_c$;
the second derivative is strictly positive. This connects the gate to
quadratic regularization \cite{ridge}, here toward a nonzero fallback.
At zero frequency the fallback is used completely; at $f_c=\kappa$ the
two scores receive equal weight; at large frequency $Y_c$ dominates.
The selected dense values are $\kappa=0.5$ and $s_0\simeq1.46055$.
With shrinkage disabled ($\kappa=0$), the procedure returns $Y_c$ directly,
including for items of zero frequency.

The frequency-based coefficient imposes a reliance policy. It does not prove
that the error variance of $Y_c$ is proportional to $1/f_c$ or that the result
is a calibrated Bayesian posterior. Also $P_c$ and $Y_c$ share sources.
Writing $Y_c=P_c+H_c$ and $\rho_c=f_c/(f_c+\kappa)$ gives
\[
 \theta(c\mid S)=\bigl[s_0+(1-s_0)\rho_c\bigr]P_c+\rho_c H_c.
\]
This is the precise separation between retained shared evidence and the
frequency-controlled additional evidence. For the selected dense recipe,
$H_c=0.5z_g+0.25z_{\rm svd}$. Rare items are not automatically boosted:
the direction of the change depends on the scores being mixed.

\section{A coherent interpretation, with distinct mathematical roles}
The resulting progression is:
\begin{enumerate}
\item Retrieve historical bundles using discounted overlap and transfer their
support to candidates. TF--IDF supplies a related frequency-weighting principle.
\item Preserve simultaneous candidate--seed--seed occurrences where selected;
these distinguish joint observations from a collection of separate pairs.
\item Share evidence when exact identity is too restrictive: compare association
profiles, pool group relations, and compare descriptions and partner descriptions.
\item Standardize the resulting scores, select their contributions on validation,
and regularize the combined estimate toward a restricted fallback using candidate
frequency.
\end{enumerate}
These connections explain what information each multiplication retains and how
it affects the score. They do not derive every fixed constant from a common
principle, establish a single joint optimization objective, or prove superiority
of the whole estimator. In particular, rarity weighting, profile compression,
mean removal, and quadratic regularization solve different mathematical problems.
Their combination is the construction to assess, rather than a claim that all
sources are versions of TF--IDF.

\begin{thebibliography}{9}
\bibitem{tfidf} Manning, Raghavan and Sch\"utze (2008).
\emph{Introduction to Information Retrieval}, inverse document frequency.
\url{https://nlp.stanford.edu/IR-book/html/htmledition/inverse-document-frequency-1.html}.
\bibitem{lsi} Manning, Raghavan and Sch\"utze (2008).
\emph{Introduction to Information Retrieval}, Chapter 18.
\url{https://nlp.stanford.edu/IR-book/pdf/18lsi.pdf}.
\bibitem{rocchio} Manning, Raghavan and Sch\"utze (2008).
\emph{Introduction to Information Retrieval}, Rocchio relevance feedback.
\url{https://nlp.stanford.edu/IR-book/html/htmledition/the-rocchio71-algorithm-1.html}.
\bibitem{rao} C. R. Rao (1982). Diversity and dissimilarity coefficients:
a unified approach. \emph{Theoretical Population Biology} 21, 24--43.
\url{https://doi.org/10.1016/0040-5809(82)90004-1}.
\bibitem{ridge} Stanford STATS 305B. Penalized regression, ridge regression.
\url{https://web.stanford.edu/class/stats305b/notes/Ridge.html}.
\end{thebibliography}
\end{document}
